\chapter{Introducción}
\thispagestyle{fancy}
\label{ch:introduccion}

En esta práctica se estudia el movimiento rectilíneo, el cual es fundamental para comprender los
principios de la mecánica clásica. Este movimiento 
se caracteriza por el desplazamiento de un objeto a lo largo de una línea recta,
y se clasifica en tres tipos principales: movimiento rectilíneo uniforme (MRU), 
movimiento rectilíneo uniformemente acelerado (MRUA)
y movimiento rectilíneo uniformemente retardado (MRUR). Cada uno de estos tipos 
de movimiento se describe mediante ecuaciones 
cinemáticas específicas que relacionan la posición, la velocidad, la 
aceleración y el tiempo.

El movimiento rectilíneo uniforme (MRU) posee una velocidad estrictamente 
constante, lo que 
implica que la aceleración es nula. En contraste, el movimiento rectilíneo 
uniformemente acelerado (MRUA) sí cambia a una tasa contante, lo 
que se traduce en una aceleración constante. Por último, el movimiento 
rectilíneo uniformemente retardado (MRUR) también presenta una aceleración 
constante, 
pero en este caso, la velocidad disminuye con el tiempo.

Gracias a esta diferencia entre sus características se pueden aplicar análisis 
específicos para cada tipo de movimiento, 
lo que permite predecir el comportamiento de los objetos en 
diferentes situaciones y estimar la incertidumbre asociada a cada resultado.

De manera complementaria, se estudia la propagación de incertidumbre, 
la cual señala que el error en una magnitud calculada a partir de otras 
magnitudes medidas directamente, 
depende del error de cada medición. En este sentido, 
el peso que cada variable tiene en la fórmula también influye en el 
resultado final. Esta práctica hace uso de dicho postulado 
con el objetivo de evaluar la 
confiabilidad de los resultados obtenidos 
frente a las limitaciones de los instrumentos de medición propuestos en el estudio 
de estos movimientos de cinemática.
